\documentclass[11pt,a4paper]{article}
\usepackage[T1]{fontenc}
\usepackage[utf8]{inputenc}
\usepackage{lmodern,amsmath,amssymb}
\usepackage[margin=25mm]{geometry}
\usepackage[hidelinks]{hyperref}
\hypersetup{pdftitle={Thermodynamics of records gives a trade-off and no floor},pdfauthor={navstokgap research working notes}}
\newcommand{\tightlist}{\setlength{\itemsep}{0pt}\setlength{\parskip}{0pt}}
\setlength{\emergencystretch}{3em}
\setcounter{secnumdepth}{0}
\begin{document}
\hypertarget{thermodynamics-of-records-gives-a-trade-off-and-no-floor}{%
\section{Thermodynamics of records gives a trade-off and no
floor}\label{thermodynamics-of-records-gives-a-trade-off-and-no-floor}}

\textbf{Result, 2026-09-27 (a recorded failure of one route).} Theorem I
of the \href{necessity-unit-and-indeterminacy.md}{unit-and-indeterminacy
note} says that back-action indeterminacy is impossible for classical
models with Liouville dynamics, product preparations obeying a prior
density ceiling, and Bayesian conditioning on records. A natural attempt
to escape it adds a heat bath and requires records to be physical:
stored, and eventually erased, at temperature \(T>0\). This note shows
that the attempt yields an exponential trade-off between dissipated
energy and posterior area, and no positive floor:

\[\eta\ \ge\ A_0\,e^{-W/(k_BT)},\]

where \(A_0\) is the prior phase-space area of the body, \(\eta\) the
area of its posterior support after the records, and \(W\) the total
work spent on measurement and erasure (Theorem 1). For every \(\eta>0\)
the construction of Theorem I reaches \(\eta\) with finitely many
finite-precision readings, hence finite \(W\). Thermodynamics therefore
prices resolution logarithmically, \(W\ge k_BT\ln(A_0/\eta)\), and
cannot single out an action. The would-be unit \(k_BT\,\tau\) built from
the bath and a record time rescales under the similarity of Theorem B of
the \href{action-unit-dimensional-selection.md}{dimensional note}. So
the thermodynamic route supplies neither ingredient (U) nor (I), and the
question of STATE item 2 stays with the three denials listed in the
unit-and-indeterminacy note.

The ingredients are standard (Landauer, Bennett, the Sagawa--Ueda bound
for measurement plus erasure); the note records their consequence for
this programme, with no novelty claimed.

\hypertarget{the-bound}{%
\subsection{1. The bound}\label{the-bound}}

Setting: the classical model of Theorem I, with body and pointers
coupled to a heat bath at temperature \(T\). Records are pointer states
that are read, used, and eventually reset to a standard state. The
generalized second law for measurement and erasure (classical form of
the bound of
\href{https://doi.org/10.1103/PhysRevLett.102.250602}{Sagawa and Ueda
2009}, metadata; the erasure part is
\href{https://doi.org/10.1147/rd.53.0183}{Landauer 1961} and
\href{https://doi.org/10.1007/BF02084158}{Bennett 1982}, metadata)
states that the total work of a measurement that gains mutual
information \(I\) (in nats) about the system, followed by erasure of the
record, satisfies \(W_{\rm meas}+W_{\rm eras}\ge k_BT\,I\).

\textbf{Theorem 1.} Let the body's prior be uniform on a phase-space
region of area \(A_0\), and let a protocol of records leave, for almost
every outcome, a body posterior supported in a region of area at most
\(\eta\). Then the total work of measurement and erasure satisfies

\[W\ \ge\ k_BT\,\ln\frac{A_0}{\eta},\qquad\text{equivalently}\qquad
\eta\ \ge\ A_0\,e^{-W/(k_BT)}.\]

\emph{Proof.} The differential entropy of the uniform prior is
\(\ln A_0\) (area measured in any fixed unit), and a density supported
in area at most \(\eta\) has differential entropy at most \(\ln\eta\).
The mutual information between body and records is the prior entropy
minus the average posterior entropy, so \(I\ge\ln(A_0/\eta)\); the unit
of area cancels. Apply the measurement-plus-erasure bound. \(\square\)

\textbf{The construction of Theorem I is thermodynamically cheap.} Its
two pointers are read to precisions \(\epsilon'\) and \(\epsilon_2'\);
each reading stores finitely many bits for a bounded pointer range, and
erasing them costs a finite multiple of \(k_BT\ln2\). So every posterior
area \(\eta>0\) is reached at finite work, and the bound of Theorem 1 is
the only thermodynamic constraint.

\hypertarget{what-this-says-about-the-necessity-question}{%
\subsection{2. What this says about the necessity
question}\label{what-this-says-about-the-necessity-question}}

\begin{itemize}
\tightlist
\item
  \textbf{No floor.} The minimal posterior area at work budget \(W\) is
  \(A_0e^{-W/(k_BT)}\), which tends to zero as the budget grows. A
  record of Newton's sagitta at resolution \(s\) with impulse resolution
  \(\delta p\) costs at least \(k_BT\ln(A_0/(s\,\delta p))\): resolving
  a phase-space area a thousand times below \(\hbar\) costs only
  \(k_BT\ln1000\) more than resolving \(\hbar\) itself. Thermodynamics
  cannot distinguish \(\hbar\).
\item
  \textbf{No unit.} The constants of the route are \(k_BT\) and whatever
  times the protocol uses. An action \(k_BT\,\tau\) rescales with
  \(\tau\), and the admitted class is closed under slowing the protocol,
  so Theorem B of the dimensional note gives floor zero. The universal
  unit of Theorem U enters through the equilibrium spectrum of
  radiation, which the record bound never uses.
\item
  \textbf{The exponential again.} The trade-off is a Boltzmann factor,
  as in the exponential ladder of the
  \href{dimension-ladder.md}{dimension-ladder note}, §4. Here the
  exponent contains no action constant, so the factor suppresses without
  setting a scale; in Wien's tail \(e^{-h\nu/(k_BT)}\) the same factor
  sets one, because \(h\) appears in the exponent.
\end{itemize}

\hypertarget{consequence-for-state}{%
\subsection{3. Consequence for STATE}\label{consequence-for-state}}

The thermodynamic route to (I) is closed at theorem level: records with
Landauer costs satisfy Theorem I with an exponential work--area
trade-off. STATE item 2 keeps its three candidate denials (of Liouville
dynamics, of product preparations, of Bayesian conditioning); the
background and inflexion routes are under evaluation in the 2026-09-27
Astra run.
\end{document}
